\begin{textsample}{NEG Dim 9 – global\_north\_2025\_03 – Score -8.00 – global\_north\_2025\_03/121965580.txt}  \label{ex:f9_neg_143}
Israel must use every tool to force Hamas compliance to release remaining hostages - editorial As so many released hostages consistently say in their brave testimonies, the captives are living in an underground hell that must end.
Tel Aviv protest calls for the release of all remaining hostages, now 512 days in \textbf{captivity}, March 1, 2025. ( photo credit : Dana Reany ) Sunday and Monday were two days in a row when thousands of people lined the streets of Tel Aviv, Rishon Lezion, and the highways and entrances flanking the southern \textbf{kibbutzim} for the \textbf{funerals} of slain hostages Shlomo Mantzur and Itzhak Elgarat, as they were laid to rest in Kissufim and Nir Oz, respectively.
The nightmarish scenario of burying people who were \textbf{kidnapped} from their homes was slightly relieved only by the show of solidarity from the thousands who stood on the sides of highways, the flags they held signaling a promise to honor, remember, and fight for a better future.
Talks continue to stall on the advancement @ @ @ @ @ @ @ @ @ @ -- which both Israel and Hamas signed for -- which is set to include the return of all the hostages, \textbf{alive} and killed, and a withdrawal of Israeli forces.
What is looking more and more likely is a return to fighting and an indefinite pause in hostage releases, unless external pressure -- i.e.
American -- succeeds in breaking the impasse.
Everyone wants the hostages home ; how to go about it is the subject of disagreement.
But, as so many released hostages consistently say in their brave testimonies, the captives are living in an underground hell that must end.
Time must be the determining factor in the next stages.
The people must keep pushing for their earliest return.
The clarity and urgency of this cause might become hazy, especially as the number of hostages drops.
It might seem like the urgency is less now because there are fewer hostages and dozens are dead.Israelis attend a rally calling for the release of hostages held by Hamas in Gaza, in Tel Aviv, March 1, @ @ @ @ @ @ @ @ @ @ is to bring them all back as soon as possible, those \textbf{alive} to their families and to begin healing, and the bodies to proper burial.
Danny Elgarat, Itzhak ' s brother, said at his \textbf{funeral} on Monday,"You are home.
You came home.
We promised, cried it out from every stage we could, that this is what we would do, that you would come home."Danny is someone who indeed fought in every way he could for his brother ' s return, as all the hostage families have done in the way they deemed right for them.
Rachel Danzig, Itzhak ' s sister, said on Monday,"You are up there with... our \textbf{parents}, looking down on us, sending us \textbf{love}.
I am down here, a survivor from \textbf{Kibbutz} Nir Oz...
We came here so many years ago and gave so much of ourselves.
We got so much in return -- until it was all taken away.
The ground turned into a field @ @ @ @ @ @ @ @ @ @.
They took you and the father of my children, Alexander Danzig, to dark, hellish \textbf{tunnels}.
You didn't return \textbf{alive}.
We failed."Israel to have US backing in any decision US President Donald Trump has already stated that he will support any decision Israel takes concerning Gaza and Hamas.
The terrorist organization has already rejected special envoy Steve Witkoff ' s proposal for a continuation of phase one of the hostage deal.
In light of Hamas ' s continued intransigence, the Israeli government made the necessary decision on Saturday night to halt humanitarian aid deliveries to Gaza.
This move is essential to exert maximum pressure on Hamas, making it clear that as long as it holds hostages and refuses a deal, it will bear the full consequences of its actions.
Though Israel has long provided humanitarian aid to Gaza -- even after October 7 and while more than 200 hostages remain captive -- this generosity has been exploited.
Innocent Israeli civilians are being held underground in dire conditions, without food, running @ @ @ @ @ @ @ @ @ @ aid for its own survival.
The time for half-measures is over.
Israel must now do everything possible to prevent a return to war while ensuring that Hamas has no choice but to release the hostages.
Every available diplomatic, economic, and military tool should be leveraged to force Hamas into compliance.
The path forward must be one of strength and resolve -- not indefinite waiting, but unyielding pressure that leads to the hostages ' swift return.
Our healing, as a nation, can only begin once all the hostages are home, and not in a process that leaves them there for an indefinite period.
Therefore, we need to be tough and demand what is just and proper -- Hamas must release our hostages, all of them.

% matched lemmas: alive, captivity, funeral, kibbutz, kidnap, love, parent, tunnel
\end{textsample}
